% Generated draft: requires mathematical and source review.
\documentclass[11pt]{article}
\usepackage[margin=25mm]{geometry}
\usepackage{fontspec}
\setmainfont{Latin Modern Roman}
\usepackage{amsmath,amssymb,mathtools,unicode-math}
\setmathfont{Latin Modern Math}
\usepackage{enumitem,needspace,xurl,hyperref}
\setlength{\parindent}{0pt}
\setlength{\parskip}{5pt}
\setlength{\emergencystretch}{3em}
\newcommand{\N}{\mathbb{N}}
\newcommand{\Z}{\mathbb{Z}}
\newcommand{\Q}{\mathbb{Q}}
\newcommand{\R}{\mathbb{R}}
\newcommand{\C}{\mathbb{C}}
\newcommand{\F}{\mathbb{F}}
\newcommand{\A}{\mathbb{A}}
\newcommand{\PP}{\mathbb{P}}
\newcommand{\E}{\mathbb{E}}
\newcommand{\eps}{\varepsilon}
\newcommand{\dd}{\,\mathrm{d}}
\begin{document}
\section*{UnsolvedMath ID: 11000277}
\textbf{Source problem code:} catalogue entry
\subsection{Definitions and mathematical statement}
Fix a free basis \(x_1,\ldots,x_n\) of \(F_n\). Let \(A_{n-1}\leq\operatorname{Aut}(F_n)\) be the embedded \(\operatorname{Aut}(\langle x_1,\ldots,x_{n-1}\rangle)\), extending each automorphism by fixing \(x_n\). The question is whether, for every group \(G\) and homomorphism \(\phi:\operatorname{Aut}(F_n)\to G\),
\[
 |\phi(A_{n-1})|<\infty\quad\Longrightarrow\quad
 |\phi(\operatorname{Aut}(F_n))|<\infty.
\]
The conclusion concerns the image, not all of \(G\), and no uniform image-size bound is requested.

The rank range is not printed. In rank \(2\), \(A_1\cong\operatorname{Aut}(\Z)\) is finite, while \(\operatorname{Aut}(F_2)\) is infinite, so the identity map shows that a restriction is necessary. The source also reports a positive block-subgroup analogue for \(\operatorname{GL}(n,\Z)\), related to algebraic \(K\)-theory; its rank hypotheses are likewise omitted.

The continuation discusses a different approach to property (T): Zuk [85] gives a sufficient link-Laplacian spectral criterion for a simply connected group complex. It suggests applying such a criterion to natural \(\operatorname{Aut}(F_n)\) or \(\operatorname{Out}(F_n)\) complexes, such as the Outer-space spine, but reports David Fisher's observation that Izeki and Natayani [55], alternatively unpublished Schoen--Wang work, obstruct that strategy. The precise complex, Laplacian normalization and obstruction hypotheses are absent. This is a limitation of the proposed method, not a claim that the groups cannot have property (T).
\subsection{Short English statement}
Does finite image on the standard rank-\(n-1\) automorphism subgroup force finite image on all of \(\operatorname{Aut}(F_n)\), in the intended rank range?
\subsection{Sources}
\begin{itemize}
\item \emph{Problems on Mapping Class Groups and Related Topics}, edited by Farb, Question 2.9 and following spectral-method discussion, supplied PDF page 331. \url{https://www.math.uchicago.edu/~farb/papers/mcgbook.pdf}. Recorded access 2026-07-29; extraction pdf-font-structure.
\item Zuk [85]; Izeki and Natayani [55] (source spelling); alternative Schoen--Wang work marked unpublished, with the obstruction observation credited to David Fisher. Complete bibliography and exact hypotheses are absent.
\item Local UnsolvedMath registry, \texttt{data/unsolvedmath/problems.json}, permanent ID \(11000277\), exact source code \texttt{catalogue entry}; supplied OPEN-TRIAGE summary says truncated.
\end{itemize}
\end{document}
